\documentclass[11pt,a4paper]{article}
\usepackage{templates/modeling-notes}
\hypersetup{pdftitle={Mathematical modeling -- lecture template}}
\title{\textbf{Mathematical Modeling}\\[2mm]\large Lecture template}
\author{Personal study notes}
\date{}
\begin{document}
\maketitle
\tableofcontents
\section{Problem and modeling idea}\label{sec:problem}
% Describe the question and the intended use of the model.
\section{Notation, assumptions, and units}\label{sec:assumptions}
% Define symbols, domains, physical units, and modeling assumptions.
\section{Definitions and results}\label{sec:results}
% Use definition, theorem, lemma, corollary, proposition, remark, proof.
\section{Derivation}\label{sec:derivation}
% Explain each step and the conditions under which it is valid.
\section{Worked example}\label{sec:example}
% Add a complete example only when studying the actual material.
\section{Limitations and common mistakes}\label{sec:limitations}
\section{Sources and editorial remarks}\label{sec:sources}
% Record source pages, dates, alignment, and editorial additions here.
% Point backwards to relevant sections; do not cite this section from the body.
% Maintain terminology in docs/glossary.md.
\end{document}
